\sekcja{Tłumaczenia\footnotemark}{}
%---------------------
\small
\vspace{-5mm}
\red{\textbf{Ant.1}} O cudowna zamiano: * Stworzyciel rodzaju ludzkiego przyjmując ciało z duszą, raczył się narodzić z Dziewicy: i stając się człowiekiem bez udziału człowieka, udzielił nam swego Bóstwa.

\red{\textbf{Psalm 109}}\\
1. Rzekł Pan Panu memu: * Siądź po prawicy mojej:\\
2. Aż położę nieprzyjaciół twoich * podnóżkiem nóg twoich.\\
3. Berło mocy twojej wypuści Pan z Syjonu: * panuj wpośród nieprzyjaciół twoich.\\
4. Przy tobie panowanie w dzień mocy twej w blaskach świętości: * z żywota przed jutrzenką zrodziłem cię.\\
5. Przysiągł Pan i nic będzie żałował: * Ty jesteś kapłanem na wieki, według porządku Melchizedecha.\\
6. Pan po prawicy twojej, * poraził w dzień gniewu swego królów.\\
7. Będzie sądził narody, dopełni zagłady, * potłucze głowy wielu na ziemi.\\
8. Ze strumienia na drodze pić będzie: * dlatego wywyższy głowę.

%---------------------
\red{\textbf{Ant.2}} Kiedyś się narodził * w niewymowny sposób z Dziewicy, wypełniły się Pisma. Zstąpiłeś jako rosa na runo, aby zbawić rodzaj ludzki! Ciebie chwalimy, Boże nasz.

\red{\textbf{Psalm 112}}\\
1. Chwalcie, słudzy, Pana, * chwalcie imię Pańskie.\\
2. Niech będzie imię Pańskie błogosławione * odtąd i aż na wieki.\\
3. Od wschodu słońca aż do zachodu, * chwalebne jest imię Pańskie.\\
4. Wywyższony jest Pan nad wszystkie narody * i nad niebiosa chwała Jego.\\
5. Któż jest jak Pan, Bóg nasz, który mieszka na wysokości, * a na niskość spogląda na niebie i na ziemi?\\
6. Podnosi z ziemi biednego * i z gnoju wywyższa ubogiego:\\
7. Aby go posadzić z książętami, * z książętami ludu swego.\\
8. Który czyni, iż niepłodna mieszka w domu, * matka synów z weselem.

%---------------------
\red{\textbf{Ant.3}} W krzaku, który widział Mojżesz, * płonącym, a nie palącym się, widzimy obraz chwalebnego dziewictwa Twego: Boża Rodzicielko, wstawiaj się za nami.

\red{\textbf{Psalm 121}}\\
1. Uweseliłem się z tego, co mi powiedziano: * Pójdziemy do domu Pańskiego.\\
2. Stały nogi nasze * w dziedzińcach twoich, Jeruzalem.\\
3. Jeruzalem, które się buduje jak miasto, * w sobie zespolone.\\
4. Bo tam wstępowały pokolenia, pokolenia Pańskie, * ustawa dla Izraela, by sławić imię Pana.\\
5. Bo tam stały stolice sądowe, * stolice dla domu Dawidowego.\\
6. Proście o to, co jest ku pokojowi dla Jeruzalem, * a dostatek niech będzie dla tych, którzy cię miłują.\\
7. Niechaj będzie pokój w warowni twojej * i dostatek w basztach twoich.\\
8. Ze względu na braci moich i bliźnich moich * wzywałem pokoju dla ciebie.\\
9. Dla domu Pana, Boga naszego, * szukałem dobra tobie.

%---------------------
\red{\textbf{Ant.4}} Zakwitła różdżka Jessego: * wyszła gwiazda z Jakóba: Dziewica porodziła nam Zbawiciela; chwalimy cię, Boże nasz.

\red{\textbf{Psalm 126}}\\
1. Jeśli Pan nie zbuduje domu, * próżno pracowali, którzy go budują.\\
2. Jeśli Pan nie będzie strzegł miasta, * próżno czuwa ten, który go strzeże.\\
3. Próżno wam przede dniem wstawać, * wstańcie, skoroście usiedli, którzy pożywacie chleba boleści.\\
4. Gdy da umiłowanym swym spanie: * oto dziedzictwem Pańskim synowie, zapłatą owoc żywota.\\
5. Jak strzały w ręku mocarza, * tak synowie utrapionych.\\
6. Błogosławiony człowiek, który napełnił nimi pragnienie swoje: * nie zawstydzi się, kiedy będzie mówił z wrogami swymi w bramie.

%---------------------
\red{\textbf{Ant.5}} Oto Marja porodziła * nam Zbawiciela, którego Jan zoczywszy, zawołał: Oto Baranek Boży, oto, który gładzi grzechy świata, alleluja.

\red{\textbf{Psalm 147}}\\
1. Chwal, Jeruzalem, Pana, * chwal, Syjonie, Boga twego.\\
2. Albowiem umocnił zawory bram twoich, * pobłogosławił synom twym w tobie.\\
3. On użyczył pokoju granicom twoim * i nasyca cię najwyborniejszym zbożem.\\
4. On puszcza słowo swe na ziemię, * bardzo prędko biegnie mowa Jego.\\
5. On daje śnieg jak wełnę, * a mgłę rozsypuje jak popiół.\\
6. Spuszcza lód swój jak bryły, * kto się ostoi przed zimnem Jego?\\
7. Pośle słowo swoje a roztopi je, * wionie wiatr Jego, a pociekną wody.\\
8. On oznajmia słowo swe Jakubowi, * ustawy i przykazania swe Izraelowi.\\
9. Nie uczynił tak żadnemu narodowi, * i nie objawił im przykazań swoich.

%---------------------
\red{\textbf{Capitulum}}\\
Oto Ja posyłam Anioła mojego, a nagotuje drogę przed obliczem mojem. A zarazem przyjdzie do kościoła swego Panujący, którego wy szukacie, i Aniół przymierza, którego wy chcecie.

%---------------------
\red{\textbf{Hymn}}\\
1. Witaj, Gwiazdo morza, * Rodzicielko Boża, * Dziewico nietknięta * Bramo niebios święta.\\
2. Gabrjel zwiastował * Tę cudu godzinę, * Obdarz nas pokojem * Mażąc Ewy winę.\\
3. Zerwij ziemskie pęta * Rzuć światłości blaski, * Niech grzeszni i ślepi * Doznają Twej łaski.\\
4. O, Matko nieś nasze * Prośby i westchnienia * Temu, co zrodzony * Dla świata zbawienia.\\
5. Dziewico przeczysta, * Prosim Cię w pokorze, * Niechaj czyste serca * Mają sługi Boże.\\
6. I bezpieczną drogą * Prowadź nas do celu, * Wtedy szczęście nasze * Znajdziem w Zbawicielu.\\
7. Chwalmy Ojca, Syna * I Ducha Świętego, * Chwalmy wiecznie Boga * W Trójcy Jedynego. * Amen.

%---------------------
\red{\textbf{Versus}}\\
\red{\Vbar.}~Odpowiedź był wziął Symeon od Ducha Świętego.\\
\red{\Rbar.}~Że nie miał oglądać śmierci, ażby pierwéj oglądał Chrystusa Pańskiego.

%---------------------
\red{\textbf{Ant.Magnificat}}~~Starzec * niósł Dziecię, lecz Dziecię starcem rządziło: Oddał cześć Temu, którego Dziewica poczęła, Dziewica zrodziła, Dziewicą zostając po narodzeniu.

\red{\textbf{Magnificat}}\\
1. Wielbij * duszo moja, Pana.\\
2. I rozradował się duch mój: * w Bogu, Zbawicielu moim.\\
3. Iż wejrzał na niskość służebnice swojéj: * albowiem oto odtąd błogosławioną mię zwać będą wszystkie narody.\\
4. Albowiem uczynił mi wielkie rzeczy, który możny jest: * i święte imię jego.\\
5. A miłosierdzie jego od narodu do narodów: * bojącym się jego.\\
6. Uczynił moc ramieniem swojem: * rozproszył pyszne myśli serca ich.\\
7. Złożył mocarze z stolice: * a podwyższył niskie.\\
8. Łaknące napełnił dobrami: * a bogacze z niczem puścił.\\
9. Przyjął Izraela, sługę swego: * wspomniawszy na miłosierdzie swoje.\\
10 Jako mówił do ojców naszych: * Abrahamowi i nasieniu jego na wieki.

%---------------------
\red{\textbf{Oratio}}\\
Wszechmogący wieczny Boże, Majestatu Twego pokornie błagamy, aby, jak Syn Twój Jednorodzony dnia dzisiejszego w postaci naszego ciała do świątyni został przyniesiony, tak abyśmy z sercem oczyszczonem przed Tobą stanąć mogli.
Przez tegoż naszego Pana Jezusa Chrystusa, Twojego Syna: który z Tobą żyje i króluje w jedności Ducha Świętego Bóg, przez wszystkie wieki wieków. Amen.
\footnotetext{Tłumaczenia pochodzą ze strony divinumofficium.com}